\documentclass[11pt]{article}

\usepackage[margin=1in]{geometry}
\usepackage[utf8]{inputenc}
\usepackage{hyperref}
\usepackage{graphicx}
\usepackage{booktabs}
\usepackage{amsmath}
\usepackage{amssymb}
\usepackage{xcolor}
\usepackage{natbib}

\hypersetup{
  colorlinks = true,
  linkcolor  = blue!50!black,
  citecolor  = blue!50!black,
  urlcolor   = blue!50!black
}

\newcommand{\acone}{\texttt{AC@1}}
\newcommand{\code}[1]{\texttt{#1}}

\title{Eval-First Agentic RCA: A Deterministic Accuracy/MTTR Regression Gate and
Shadow-to-Active Rollout for Shipping LLM-in-the-Loop Reliability Tooling}

\author{Victor Robin\\ BlueRobin (independent homelab research)}

\date{June 2026}

\begin{document}

\maketitle

\begin{abstract}
We describe the evaluation methodology that underpins the BlueRobin Debug Agent, an
LLM-in-the-loop root-cause-analysis (RCA) service that runs on a self-hosted homelab under
a hard ${\sim}50$~EUR/month budget. The methodology is \emph{eval-first}: before any
reasoning component (graph blame-propagation, change-point detection, hypothesis-state
machine, episodic memory) is allowed to govern production behaviour, it must pass a
deterministic offline replay harness that scores the \emph{real} RCA path against a labelled
ground-truth incident set and emits an Accuracy@1 (\acone{}) plus
mean-time-to-resolution (MTTR/TTM) scorecard. A per-commit \textbf{regression gate} with
explicit tolerances (\acone{} drop $\le 0.02$, MTTR rise $\le 30$~s) blocks any merge that
degrades localisation accuracy or inflates analysis latency. The harness has two arms: a
\emph{scorer-only} arm where the labelled cause is planted as the unambiguous top-1 (a
regression sentinel, $\acone{} = 1.0$ by construction) and an honest \emph{live-path} arm
that injects \textbf{distractor candidates} with higher topology criticality but lower
error-rate, so $\acone{} < 1.0$ is achievable on the production ranking path. New components
ship behind a \textbf{shadow-to-active} rollout (\code{RcaLiveActivation.Mode}~$\in
\{\text{Shadow}, \text{Active}\}$): they compute and log their would-be decisions for an
observation period before they are permitted to govern. We are explicit that the
ground-truth set is small (9 self-seeded incidents spanning $\ge 3$ fault classes), that
MTTR figures are sub-second replay-clock proxies rather than production MTTR, and that the
authors themselves flag the by-construction $\acone{} = 1.0$ as an ``oversimplified
benchmark.'' The contribution is not a new accuracy record; it is a reproducible discipline
--- change-correlation as the primary signal, a deterministic merge gate, and staged
activation --- that makes the rest of the system measurable and safe to ship.
\end{abstract}

\section{Introduction}

Automated RCA for microservice and Kubernetes systems has converged on LLM agents that
interleave reasoning and tool calls (ReAct; \citealp{yao2022react}) over observability
backends. The appeal is obvious: an agent can query traces, logs, metrics, deploy history
and source code, then narrate a hypothesis. The problem is equally obvious once one tries
to ship it: \textbf{an LLM-only RCA loop has no native notion of being right, and no native
notion of regressing.} A prompt tweak, a model swap, or a new tool can silently degrade
localisation accuracy, and an operator would never know, because the only feedback signal
is a human reading a plausible-sounding paragraph during an incident.

The independent literature is blunt about the ceiling. The best agent on OpenRCA reaches
${\sim}11.3\%$ exact-match \citep{xu2025openrca}; the best frontier model on IBM's
ITBench-AA reaches ${\sim}47\%$ precision@full-recall \citep{aa2026itbenchaa}. Benchmarks
themselves are now suspected of overstating capability: simple rule-based methods match
``SOTA'' on several of them \citep{rethinking2025faultprop}. In this environment, the
responsible engineering posture is \textbf{advisory / human-in-the-loop} --- a human always
merges the fix --- and the responsible \emph{development} posture is \textbf{eval-first}:
the eval comes before the feature, and every commit is a regression test (HolmesGPT runs
150+ evals per commit; \citealp{holmesgpt,evaldriven}).

This paper documents how the BlueRobin Debug Agent operationalises that posture. The system
is real, implemented in .NET~10 over FalkorDB and Qdrant, and governed by requirement
\code{FEAT-027} (Stack-Graph RCA) extending \code{FEAT-008} (Observability \& Auto-Debug).
Our contributions:

\begin{itemize}
  \item \textbf{A deterministic offline replay harness} (\code{RcaEvalHarness}) that
  replays a labelled ground-truth incident set through the \emph{actual} production RCA path
  and emits a per-incident and aggregate \acone{} + MTTR scorecard, byte-for-byte
  reproducible (\code{NFR-SG-5}).
  \item \textbf{A per-commit regression gate} (\code{RcaEvalRegressionGate}) with explicit,
  committed tolerances (\acone{} drop $\le 0.02$, MTTR rise $\le 30$~s) wired into CI as a
  merge blocker.
  \item \textbf{An honest live-path arm} that injects distractor candidates --- siblings
  with higher topology criticality but lower error-rate --- so that $\acone{} = 1.0$ is
  \emph{not} guaranteed and the ranker's accuracy on the production path is actually
  measured.
  \item \textbf{A shadow-to-active rollout discipline} (\code{RcaLiveActivation}, ADR-039)
  under which the verification gate, BARO scorer, evidence soft-cap and incident-memory
  recall each run in \emph{shadow} (compute + log, do not govern) before promotion to
  \emph{active}.
  \item \textbf{Change-correlation as the primary ranking signal}, motivated by the Google
  SRE finding that ${\sim}70\%$ of outages follow a change to a live system
  \citep{beyer2016sre}.
  \item \textbf{An honest threats-to-validity treatment}: we state the small-N, self-seeded,
  replay-clock limitations plainly rather than burying them.
\end{itemize}

\section{Related Work}

\paragraph{Eval-driven development for agents.}
The discipline of treating evals as per-commit regression tests, not occasional report
cards, is crystallised by the HolmesGPT project (Robusta + Microsoft; CNCF Sandbox), which
runs 150+ evals on every commit, and by the broader ``the eval comes first'' movement
\citep{holmesgpt,evaldriven}. Datadog's Bits AI SRE similarly scores against a labelled set
with an LLM-judge and surfaces an agent trace \citep{datadog2026bits}. Our harness adopts
this discipline but constrains the per-commit oracle to be \textbf{deterministic and
LLM-free} so it can gate merges cheaply and reproducibly; the LLM-judge is relegated to an
optional, scheduled, token-capped path.

\paragraph{Benchmarks for RCA.}
OpenRCA \citep{xu2025openrca} provides the most honest ``how bad are LLMs at RCA'' number
(${\sim}11.3\%$ exact-match over 335 cases). AIOpsLab \citep{chen2025aiopslab} offers 48
problems across four levels with TTD/TTM, step-count and token-cost metrics over live
Kubernetes (DeathStarBench) with ChaosMesh faults. ITBench \citep{jha2025itbench}
introduces a Normalised Topology-Aware Metric and reports MTTD 72~s / MTTR 282~s, and shows
traces are load-bearing (removing them drops diagnosis $13.8\% \rightarrow 9.5\%$). RCAEval
\citep{pham2024rcaeval} is the cleanest reproducible baseline set --- 735 cases, 11 fault
types, 15 baselines --- and standardises the \acone{}, \code{Avg@k}, \code{MRR} metric
family we adopt. Crucially, the Fault-Propagation-Aware critique
\citep{rethinking2025faultprop} argues current RCA benchmarks are oversimplified and that
published scores overstate real capability. We treat all of these as \textbf{external
sanity yardsticks only}, never as dependencies, and we adopt their critique as our own
Section~\ref{sec:limitations} caveat.

\paragraph{Change / deploy correlation.}
The single most predictive RCA signal is recent change. The Google SRE Book states that
${\sim}70\%$ of outages are due to changes in a live system \citep{beyer2016sre}. The
Microsoft Teams incident study \citep{ghosh2022teams} decomposes root-cause categories ---
Code Bug 27\%, Dependency 16.4\%, Infra 15.8\%, Deployment 13.2\%, Config 12.5\% --- and
reports that ${\sim}80\%$ of incidents are mitigated \emph{without} a code/config fix, with
rollback the \#1 mitigation. The classic configuration-error study \citep{yin2011config}
found 16.1--47.3\% of misconfigurations cause full unavailability or severe degradation.
These motivate first-class \code{Deploy}/\code{Change} graph nodes and an onset-window
change term in our ranker (\code{FR-SG-25}, \code{FR-SG-26}).

\paragraph{Agent failure modes.}
The MAST failure taxonomy \citep{cemri2025mast} (310 ITBench traces) finds that incorrect
self-verification, premature termination, memory loss, and reasoning-action mismatch
account for ${\sim}94\%$ of agent failures, with the prescriptive lessons: never let the LLM
grade its own homework; require hard tool evidence before exit; enforce termination via an
FSM outside the model. Our gate and shadow-to-active discipline are direct responses; the
verification-gate FSM and hypothesis-state machine are detailed in companion papers in this
series.

\section{Method / Architecture}
\label{sec:method}

\subsection{The deterministic replay harness}
\label{sec:harness}

The canonical eval is \code{RcaEvalHarness}
(\code{src/DebugAgent/StackGraph/Eval/RcaEvalHarness.cs}), an offline .NET test/CLI
entrypoint --- \textbf{not} a new in-cluster service --- that replays the labelled
ground-truth set through the real RCA path (\code{StackGraphRcaService} $\rightarrow$
\code{RootCauseScorer}/\code{OnsetWindowChangeScorer} over a seeded FalkorDB) and emits a
deterministic \acone{} + MTTR/TTM scorecard, per-incident and aggregate. It is the
source-of-truth realisation of \code{FR-SG-23} (ADR-036~D1).

The data flow is: \code{ground-truth-incidents.json} (9 labelled incidents) is read by the
harness, which seeds a clean FalkorDB graph (\code{bluerobin\_stack\_eval}), runs the real
\code{StackGraphRcaService} ranking path, and produces an \code{EvalScorecard}
(\acone{} + MTTR per incident). That scorecard is fed, together with the committed
\code{scorecard-baseline.json}, to the \code{RcaEvalRegressionGate}: within tolerance, the
merge is allowed; on regression, the gate exits 1 and the merge is blocked.

\paragraph{Two arms.}
The harness deliberately exposes two scoring paths:

\begin{itemize}
  \item \code{ReplayAsync} --- the \textbf{scorer-only} arm. Each incident's labelled root
  cause is seeded as the unambiguous top-1 candidate (\code{RootCauseCriticality = 0.9},
  \code{RootCauseErrorRate = 0.9}), so $\acone{} = 1.0$ \emph{by construction}. This is a
  \textbf{regression sentinel}, not an accuracy measurement: it proves that a code change
  did not break the plumbing that maps a seeded subgraph to the correct top-1.
  \item \code{ReplayLivePathAsync} --- the \textbf{live-path} arm. It runs onset-from-alert
  and injects \textbf{distractor candidates}: sibling services the symptom also calls, with
  \code{DistractorErrorRate = 0.55} (non-trivial but below the root cause's 0.9) and
  \code{DistractorCriticality = 0.95} (deliberately \emph{higher} topology than the root
  cause's 0.9 --- a topology trap). Because a distractor outranks the true cause on the
  topology term, $\acone{} < 1.0$ is now achievable. This is the honest production-path
  accuracy gate.
\end{itemize}

The distractor design is the methodological crux. Without it, the eval is gameable: any
ranker that trusts topology blindly still scores 1.0. With it, a candidate that boosts
topology weight at the expense of the error-rate / change-onset signal will mis-localise GT
incidents and drop \acone{}, failing the gate. The trap is exactly the failure mode the
literature warns about --- topology criticality is necessary but not sufficient; the
discriminating signal is the per-service error-rate plus the in-window change.

\subsection{The ground-truth incident set}
\label{sec:groundtruth}

The labelled set lives in \code{src/DebugAgent/eval/ground-truth-incidents.json}
(\code{schemaVersion: 1}, \code{windowSeconds: 1800}). It is \textbf{seeded only from
existing faults --- no ChaosMesh} (\code{RISK-1}): the \code{bluerobin-e2e} synthetic-test
family, the \code{bluerobin-infra} synthetic-registry scenarios (the OBS-9 login break is
canonical), and past \code{agent\_runs} incidents replayed offline. There are 9 labelled
incidents spanning $\ge 3$ fault classes (auth/login, data/datastore, deploy/config) so the
set is not trivially gameable by a single-class heuristic.

\begin{table}[h]
\centering
\caption{The labelled ground-truth incident set (9 incidents, $\ge 3$ fault classes).}
\label{tab:groundtruth}
\begin{tabular}{llll}
\toprule
ID & faultClass & expectedRootCause & changeInduced \\
\midrule
GT-01-login-break          & auth/login    & authelia              & false \\
GT-02-rca-gateway-down-floor & deploy/config & mcp-gateway         & false \\
GT-03-bad-deploy-archives-api & deploy/config & archives-api       & \textbf{true} \\
GT-04-falkordb-down        & data/datastore & falkordb             & false \\
GT-05-postgres-meshed-reset & data/datastore & postgres            & false \\
GT-06-upload-thumbnail-fail & data/datastore & archives-worker-infra & false \\
GT-07-bad-commit-web       & deploy/config & bluerobin-web         & \textbf{true} \\
GT-08-rag-chat-down        & auth/login    & authelia              & false \\
GT-09-ollama-embeddings-down & data/datastore & ollama             & false \\
\bottomrule
\end{tabular}
\end{table}

Two incidents (GT-03, GT-07) are explicitly \textbf{change-induced} --- each carries a
\code{Deploy}/\code{Change} node inside the 30-minute onset window. They are the regression
test for change-correlation: every active-mode ablation must keep $\acone{} = 1.0$ on GT-03
and GT-07 (no true-positive suppression of the change signal). The same JSON is the single
source for both the eval set and the simulation/chaos catalog (\code{OQ-CATALOG-SHARE}), so
the labelled truth and the on-demand simulation harness can never drift apart.

\subsection{Metrics and the oracle}
\label{sec:metrics}

\acone{} is an exact-string oracle: the produced top-1 candidate's natural key must equal
the labelled \code{ExpectedRootCauseNaturalKey}. There is \textbf{no LLM-judge and no human
in the per-commit path}. MTTR/TTM is the replay wall-clock from onset-replay-start to the
pipeline returning its conclusion (a \code{Stopwatch} around \code{InvestigateAsync}) --- a
time-to-mitigate \emph{proxy}, not production MTTR. With a single labelled cause per incident
scored at top-1, \code{Avg@k} and \code{MRR} reduce to \acone{} here; they are computed but
carry no extra information until multi-candidate labels land (ADR-036~D1).

\begin{table}[h]
\centering
\caption{Metrics computed by the harness and their gating role.}
\label{tab:metrics}
\begin{tabular}{lll}
\toprule
Metric & Definition (this harness) & Role \\
\midrule
\acone{}        & $1$ iff top-1 candidate key $=$ labelled root-cause key & Primary; gated first \\
MTTR/TTM        & replay wall-clock onset $\rightarrow$ conclusion        & Secondary; gated, 30~s tolerance \\
\code{Avg@k}, \code{MRR} & reduce to \acone{} (single-label)              & Reserved for multi-label future \\
\bottomrule
\end{tabular}
\end{table}

\subsection{The regression gate and tolerances}
\label{sec:gate}

The gate is \code{RcaEvalRegressionGate.Evaluate}, run by \code{RcaEvalGateRunner} with
committed default tolerances:

\begin{quote}
\begin{verbatim}
public static readonly GateTolerance DefaultTolerance =
    new(AcAt1Drop: 0.02, MttrSecondsIncrease: 30.0);
\end{verbatim}
\end{quote}

\acone{} (higher-is-better) is checked \textbf{first}: a candidate scorecard whose aggregate
\acone{} drops more than 0.02 below the committed baseline fails with exit code~1 and the
regressed metric named ``AC@1''. MTTR (lower-is-better) is checked second: a rise of more
than 30~s fails as ``MTTR''. A drift \emph{within} tolerance passes --- the 30~s MTTR band
deliberately absorbs machine-dependent replay-clock jitter (sub-second MTTR values are
explicitly not asserted exactly). The committed baseline \code{\_comment} states the
contract verbatim: it ``FAILS any candidate that regresses \acone{} below 0.98 or inflates
MTTR by $>30$~s.''

This gate is the load-bearing artefact of the whole programme: it is what makes every
\emph{later} claim in this paper series --- graph blame-propagation, BARO onset detection,
evidence-chain confidence, episodic memory --- \emph{measurable}. None of them can be merged
if they regress localisation on the ground-truth set.

\subsection{Change correlation as the primary signal}
\label{sec:change}

Because ${\sim}70\%$ of outages follow a change \citep{beyer2016sre}, the ranker treats
in-window change as a first-class term. The pre-LLM scoring blend is a re-normalised
additive blend over present terms. The base 4-term blend (\code{OnsetWindowChangeScorer})
and the 5-term BARO blend (\code{BaroBlendScorer}) are shown in Table~\ref{tab:blend}.

\begin{table}[h]
\centering
\caption{Pre-LLM scoring-blend weights. $w_{\text{topo}}$ topology criticality,
$w_{\text{tele}}$ telemetry/error-rate, $w_{\text{hist}}$ historical prior,
$w_{\text{chg}}$ in-window change, $w_{\text{baro}}$ BARO onset.}
\label{tab:blend}
\begin{tabular}{lcccccc}
\toprule
Blend & $w_{\text{topo}}$ & $w_{\text{tele}}$ & $w_{\text{hist}}$ & $w_{\text{chg}}$ & $w_{\text{baro}}$ & $\Sigma$ \\
\midrule
4-term onset (\code{OnsetWindowChangeScorer}) & 0.40 & 0.30 & 0.15 & 0.15 & --- & 1.00 \\
5-term BARO (\code{BaroBlendScorer})          & 0.35 & 0.25 & 0.15 & 0.15 & 0.10 & 1.00 \\
Fail-open floor (no chg/baro)                 & 0.50 & 0.35 & 0.15 & --- & --- & 1.00 \\
\bottomrule
\end{tabular}
\end{table}

A candidate is boosted by $w_{\text{chg}}$ when its owning service had a
\code{Deploy}/\code{Change} node in the 30-minute onset window ($[T-W, T]$, $W = 1800$~s).
When the change or BARO term is absent, the blend re-normalises over the present terms (the
fail-open floor), and that collapse is byte-for-byte identical to the prior baseline --- so a
missing signal can never \emph{introduce} a regression. The distractor arm
(Section~\ref{sec:harness}) is precisely what forces this term to earn its weight: a
distractor with higher topology but no in-window change must not outrank a true
change-induced cause, or GT-03/GT-07 fail.

\subsection{Shadow-to-active rollout}
\label{sec:shadow}

New reasoning components do not light up on merge. They ship behind
\code{RcaLiveActivation.Mode} (\code{src/DebugAgent/Rca/RcaLiveActivation.cs}, ADR-039~D1):

\begin{quote}
\begin{verbatim}
public enum RcaLiveActivationMode
{
    Shadow = 0, // default(enum) — compute + log "would-do", never govern
    Active,     // the gate / BARO / memory drive the real decision
}
\end{verbatim}
\end{quote}

In \code{Shadow} (the default), the verification gate and BARO scorer compute and log what
they \emph{would} do, but the legacy path still governs; the eval collapses byte-for-byte to
the prior baseline. Promotion to \code{Active} is an explicit operator config change
(ConfigMap-backed), not a code default. This gives a clean ablation structure: for each
increment we assert (a) \textbf{active}: no regression versus baseline; (b) \textbf{shadow}:
collapses to baseline exactly (the ``feature off'' control); (c) \textbf{onset
true-positives preserved}: GT-03/GT-07 stay at $\acone{} = 1.0$. The same shadow-first flag
governs the evidence-chain soft-cap (\code{RcaEvidenceCap.Mode}) and incident-memory recall
(\code{RcaMemory.Mode}).

This is the methodology that makes the system \textbf{safe to ship}: a new component is
observed against the gate in shadow before it is ever allowed to change an operator-facing
conclusion.

\section{Implementation}
\label{sec:impl}

The harness is pure .NET~10. Determinism (\code{NFR-SG-5}) is engineered, not hoped for:

\begin{itemize}
  \item A clean graph per run (\code{MATCH (n) DETACH DELETE n}), no live-graph reads, empty
  history priors seeded.
  \item Stable incident and distractor ordering (JSON order), so the per-commit gate is not
  flaky.
  \item An isolated FalkorDB key, \code{bluerobin\_stack\_eval} --- never the live
  \code{bluerobin\_stack} graph.
  \item Fixed onset epochs (\code{1700000000+}) so scoring has no wall-clock dependence; only
  MTTR is wall-clock, and it is absorbed by the 30~s tolerance.
  \item A reproducibility assertion: two consecutive runs must agree on aggregate \acone{} to
  $10^{-9}$ and produce identical incident ordering.
\end{itemize}

\paragraph{CI wiring.}
\code{debug-agent-pipeline.yml} runs the deterministic gate per-commit on every push to
main: it spins up \code{falkordb/falkordb:v4.4.1} on an isolated port, runs
\code{dotnet run \ldots{} -- --job=eval} against \code{ground-truth-incidents.json} and
\code{scorecard-baseline.json}, and fails the build on regression. StackGraph and lifecycle
integration suites run via Testcontainers against real FalkorDB and \code{postgres:16} ---
there is no DB mock. The optional LLM-judge (\code{RcaEvalJudgeOptions}, default OFF in CI,
hard \code{PerRunTokenCap: 50\_000}) runs only on the nightly \code{eval-nightly.yml} cron
and \textbf{never gates merges} --- this is what keeps the judge inside the
20~EUR/month LLM ceiling.

This work is governed by \code{FEAT-027}, specifically \code{FR-SG-23} (the eval harness +
\acone{}/MTTR scorecard), \code{FR-SG-24} (the regression gate + bounded LLM-judge),
\code{FR-SG-25} (\code{Deploy}/\code{Change} nodes), and \code{FR-SG-27} (Grafana deploy
annotations), under \code{NFR-SG-5} (deterministic reproducibility) and ADR-036 (eval
harness + change correlation) / ADR-039 (Phase-B shadow$\rightarrow$active live activation).

\paragraph{The ${\sim}50$~EUR/month angle.}
The entire programme runs inside a hard platform ceiling of ${\sim}50$~EUR/month: a single
Hetzner CX43 worker (${\sim}17.5$~EUR/mo), Cloudflare R2 and Tunnel on free tiers, local
Ollama embeddings, and a hard 20~EUR/month LLM token ceiling (\code{COST-7}, with
fail-closed at 100\% budget). Two design choices fall directly out of that ceiling. First,
the per-commit oracle is deterministic and LLM-free, so the gate that runs on \emph{every}
commit costs zero tokens. Second, the eval reuses the \emph{already-deployed} FalkorDB
instance under an isolated key rather than standing up a benchmarking cluster --- a live
ChaosMesh-style fault-injection benchmark over a separate Kubernetes environment is simply
not affordable, which is why the ground-truth set is seeded from existing faults. The cost
constraint is not a footnote; it is the reason the methodology looks the way it does.

\section{Evaluation}
\label{sec:eval}

The committed scorecards are the baselines the gate protects. The \textbf{scorer-only}
baseline (\code{eval/scorecard-baseline.json}) is $\acone{} = 1.0$ by construction across
all 9 incidents --- a regression sentinel. The \textbf{live-path} baseline
(\code{src/DebugAgent/eval/scorecard-live-path-baseline.json}) is the honest production-path
number: even with distractors injected, all 9 incidents still localise at top-1, so the
measured live-path \acone{} is also 1.0 on this set.

\begin{table}[h]
\centering
\caption{Committed baseline scorecards (the baselines the per-commit gate protects).}
\label{tab:scorecards}
\begin{tabular}{lcccccl}
\toprule
Arm & \acone{} & MTTR (s) & Incidents & \code{Avg@k} & \code{MRR} & Interpretation \\
\midrule
Scorer-only (\code{ReplayAsync})        & 1.0 & 0.0017 & 9 & 1.0 & 1.0 & By construction; sentinel \\
Live-path (\code{ReplayLivePathAsync})  & 1.0 & 0.0032 & 9 & 1.0 & 1.0 & Honest path; distractors \\
\bottomrule
\end{tabular}
\end{table}

The scorer-only arm is a regression sentinel, \textbf{not} an accuracy claim. The live-path
arm is the honest production path: distractors are injected, and $\acone{} < 1.0$ is
achievable. Per-incident MTTR ranges 0.0007--0.0073~s (scorer-only) and 0.0016--0.0126~s
(live-path) --- sub-second, machine-dependent, and \textbf{not asserted exactly}.

\paragraph{How to read these numbers honestly.}
The live-path arm proves the ranker is not fooled by the topology trap on \emph{these} nine
incidents --- a real result, because $\acone{} = 1.0$ there is \emph{not} guaranteed by the
harness; the distractor with higher criticality could legitimately win. But nine
self-seeded incidents is a tiny, low-diversity set, and a perfect score on it says far more
about the absence of regressions than about absolute capability. For external context: the
best agent on OpenRCA reaches ${\sim}11.3\%$ exact-match \citep{xu2025openrca} and the best
frontier model on ITBench-AA reaches ${\sim}47\%$ precision@full-recall
\citep{aa2026itbenchaa} on far harder, larger, adversarial sets. Our 1.0 is \textbf{not}
comparable to those numbers and we do not claim it is. The harness's own source calls the
by-construction arm ``the \S12 oversimplified-benchmark caveat'' --- an in-code
acknowledgement that the planted 1.0 is plumbing, not accuracy.

What the evaluation \emph{does} establish is the property we actually need: every increment
(verification gate + BARO, evidence soft-cap, memory recall) passes the gate in both active
and shadow modes, and the change-induced incidents GT-03/GT-07 retain $\acone{} = 1.0$ under
every active-mode ablation. That is a no-regression guarantee, not a capability record ---
which is exactly the claim an eval-first methodology should make.

\section{Limitations \& Threats to Validity}
\label{sec:limitations}

We state these plainly; they are not hedges, they are the honest scope of the work.

\begin{itemize}
  \item \textbf{Small, self-seeded ground truth (construct + external validity).} Nine
  incidents, seeded by the authors from existing faults, with no ChaosMesh or live fault
  injection. The set spans $\ge 3$ fault classes specifically to avoid single-class gaming,
  but it is small and low-diversity. $\acone{} = 1.0$ on it does not generalise to
  production RCA difficulty.
  \item \textbf{Author-labelled oracle (internal validity).} The expected root cause for each
  incident is an author judgement. Exact-string matching is unambiguous \emph{given} the
  label, but the labels themselves are not independently adjudicated.
  \item \textbf{MTTR is a replay-clock proxy (construct validity).} Reported MTTR is
  sub-second offline wall-clock around \code{InvestigateAsync}, not production
  fire-to-resolved MTTR. Real lifecycle MTTD/MTTA/MTTR p50/p95 are computed separately from
  live data (\code{GetKpisEndpoint}, \code{SLO-15}/\code{SLO-16}) but no measured production
  values are committed here.
  \item \textbf{By-construction $\acone{} = 1.0$ (presentation risk).} The scorer-only arm is
  a sentinel; reported naively it overstates capability. We mitigate by always reporting the
  live-path arm alongside it and by labelling the scorer arm as plumbing.
  \item \textbf{No external-baseline comparison.} OpenRCA / ITBench-AA / AIOpsLab are
  external yardsticks, never dependencies; we do not run them, so our numbers are not
  cross-comparable. This is deliberate (cost and reproducibility) but it means we cannot
  claim relative standing against published systems.
  \item \textbf{Benchmark-oversimplification critique applies to us too.} The
  Fault-Propagation-Aware critique \citep{rethinking2025faultprop} --- that simple methods
  match SOTA on easy benchmarks --- applies with full force to a 9-incident self-seeded set.
  We adopt the critique rather than resist it.
  \item \textbf{Advisory-only by design.} A human always merges; nothing auto-remediates
  (ADR-030). We position this as matching the state of the art given the ${\sim}11$--$47\%$
  autonomous ceiling \citep{xu2025openrca,aa2026itbenchaa}, not as a temporary limitation.
\end{itemize}

\section{Conclusion}
\label{sec:conclusion}

LLM-only RCA loops have no native sense of correctness or regression, which makes them
dangerous to evolve and hard to ship responsibly. The BlueRobin Debug Agent addresses this
not with a bigger model but with a discipline: a deterministic offline replay harness over a
labelled ground-truth set, a per-commit regression gate with explicit \acone{}/MTTR
tolerances, an honest distractor-injected live-path arm that makes $\acone{} < 1.0$
achievable, change-correlation as the primary signal, and a shadow-to-active rollout that
observes every new reasoning component before it is allowed to govern. None of this produces
a record accuracy number --- by design, and we say so. What it produces is the precondition
for everything else in this series: a way to add graph reasoning, change-point detection,
evidence-grounded confidence and episodic memory to an LLM-in-the-loop RCA agent and
\emph{know}, on every commit, that one did not make it worse --- all inside a
${\sim}50$~EUR/month homelab budget. The honest contribution is the gate, not the score.

\bibliographystyle{plainnat}
\bibliography{references}

\end{document}
